\chapter{Methods}
\label{ch:methods}

\section{DeepEF Architecture Overview}

The Protein Energy Model (PEM) predicts per-residue energy contributions that sum to a protein's total folding energy. The $\Delta\Delta G$ for a mutation is computed as:
\begin{equation}
    \Delta\Delta G = E_{\text{mutant}} - E_{\text{wild-type}} = (E_{\text{unfolded}}^{\text{mut}} - E_{\text{folded}}^{\text{mut}}) - (E_{\text{unfolded}}^{\text{wt}} - E_{\text{folded}}^{\text{wt}})
\end{equation}

\subsection{Input Representation}

For each protein, the input feature vector per residue is:
\begin{equation}
    \mathbf{F}_i = [\mathbf{D}_i \,||\, \mathbf{F}_b \,||\, \mathbf{e}_i \,||\, \mathbf{h}_i] \in \mathbb{R}^{16 + 32 + E + 20}
\end{equation}
where:
\begin{itemize}
    \item $\mathbf{D}_i \in \mathbb{R}^{16}$: Distance matrix features (16 Gaussian kernels applied to atom pair distances, summed over neighbors, L2-normalized)
    \item $\mathbf{F}_b \in \mathbb{R}^{32}$: Bonded features (16 bonded + 16 non-bonded Gaussian distance kernels from immediate sequential neighbors)
    \item $\mathbf{e}_i \in \mathbb{R}^{E}$: Protein language model embedding ($E = 1024$ for ProtT5, $E = 1280$ for SaProt/ESM-2)
    \item $\mathbf{h}_i \in \mathbb{R}^{20}$: One-hot amino acid encoding
\end{itemize}

The distance matrix uses Gaussian radial basis functions:
\begin{equation}
    D_{ij} = \text{ReLU}\left(\exp(\gamma \cdot d_{ij}^2)\right), \quad \gamma = -0.08
\end{equation}
computed between backbone atom pairs (N, CA, C, CB) using 16 distance bins.

\subsection{Dual-Branch GNN}

The model processes protein graphs through two parallel branches:

\textbf{GCN Branch} (bonded + non-bonded context):
\begin{itemize}
    \item Input: $[\mathbf{D}(16) \,||\, \mathbf{F}_b(32) \,||\, \mathbf{h}(20)] = 68$ dimensions (or $52$ without projection)
    \item 3 GCN layers with dropout ($p = 0.2$)
    \item Edge connectivity: fully connected (all residues connected)
    \item Output: 36-dimensional node features
\end{itemize}

\textbf{GAT Branch} (non-bonded attention):
\begin{itemize}
    \item Input: $[\mathbf{D}(16) \,||\, \mathbf{h}(20)] = 36$ dimensions
    \item 3 GATv2 layers with 4 attention heads, dropout ($p = 0.2$)
    \item Edge connectivity: $k$-NN graph ($k = 30$) based on CA atom distances
    \item Output: 36-dimensional node features
\end{itemize}

\subsection{$k$-NN Graph Construction}

For the GAT branch, we construct a directed $k$-nearest neighbor graph:
\begin{equation}
    \mathcal{E} = \{(i, j) : j \in \text{kNN}(i, k=30) \text{ based on } \|\mathbf{r}_i^{CA} - \mathbf{r}_j^{CA}\|_2\}
\end{equation}

This replaces the original fully-connected graph and provides:
\begin{itemize}
    \item Reduced computation: $O(Nk)$ vs $O(N^2)$ edges
    \item Focus on local structural contacts (which determine stability)
    \item Empirically validated: +0.043 PCC improvement over fully-connected
\end{itemize}

\subsection{Feature Fusion and Prediction}

After both GNN branches:
\begin{enumerate}
    \item Concatenate GCN and GAT outputs: $\mathbf{x} = [\mathbf{x}_{\text{GCN}} \,||\, \mathbf{x}_{\text{GAT}}] \in \mathbb{R}^{72}$
    \item Apply instance normalization
    \item Concatenate raw LLM embeddings: $\mathbf{x} = [\mathbf{x} \,||\, \mathbf{e}] \in \mathbb{R}^{72 + E}$
    \item Apply Light Attention mechanism
    \item Two-layer MLP: $(72 + E) \rightarrow 128 \rightarrow 1$
\end{enumerate}

\subsection{Light Attention}

Light Attention \cite{starklight} provides a learnable pooling mechanism over the sequence dimension:
\begin{equation}
    \mathbf{z} = \text{softmax}(\mathbf{W}_a \mathbf{X}^T) \cdot \mathbf{X}
\end{equation}
where $\mathbf{X} \in \mathbb{R}^{N \times d}$ is the sequence of node features and $\mathbf{W}_a$ learns which residues to attend to for the global prediction.

\section{Loss Function}

We use a composite loss combining Huber loss and ranking loss:
\begin{equation}
    \mathcal{L} = \mathcal{L}_{\text{Huber}} + \lambda \cdot \mathcal{L}_{\text{ranking}}
\end{equation}

\textbf{Huber Loss} ($\delta = 1.0$):
\begin{equation}
    \mathcal{L}_{\text{Huber}}(y, \hat{y}) = \begin{cases}
        \frac{1}{2}(y - \hat{y})^2 & \text{if } |y - \hat{y}| \leq \delta \\
        \delta(|y - \hat{y}| - \frac{1}{2}\delta) & \text{otherwise}
    \end{cases}
\end{equation}

\textbf{Ranking Loss} (margin = 0.1, $\lambda = 0.1$):
\begin{equation}
    \mathcal{L}_{\text{ranking}} = \frac{1}{|P|} \sum_{(i,j) \in P} \max(0, -\text{sign}(y_i - y_j)(\hat{y}_i - \hat{y}_j) + m)
\end{equation}
where $P$ is the set of mutation pairs within a protein, encouraging correct relative ordering.

\section{Training Pipeline}

\subsection{Training Configuration}

\begin{table}[h]
\centering
\caption{Training hyperparameters}
\label{tab:hyperparams}
\begin{tabular}{lc}
\toprule
\textbf{Parameter} & \textbf{Value} \\
\midrule
GNN layers & 3 \\
Dropout rate & 0.2 \\
Gaussian coefficient $\gamma$ & --0.08 \\
Batch size & 1 (protein-level) \\
Mini-batch size & 64 (mutations per step) \\
Learning rate & $10^{-4}$ \\
LR schedule & Cosine annealing ($\eta_{\min} = 10^{-6}$) \\
Optimizer & Adam (weight decay $= 10^{-5}$) \\
Gradient accumulation & 4 steps \\
Epochs & 30 \\
$\Delta\Delta G$ clipping & $[-1.0, 5.0]$ kcal/mol \\
\bottomrule
\end{tabular}
\end{table}

\subsection{$\Delta\Delta G$ Computation}

For each protein in a batch:
\begin{enumerate}
    \item Construct folded graph from backbone coordinates + embedding of mutation $m$
    \item Construct unfolded graph (diagonal-only distance matrix, same features)
    \item Predict energies: $E_{\text{folded}}(m)$, $E_{\text{unfolded}}(m)$
    \item $\Delta G(m) = E_{\text{unfolded}}(m) - E_{\text{folded}}(m)$
    \item $\Delta\Delta G(m) = \Delta G(m) - \Delta G(\text{wild-type})$ (index 0 is always WT)
\end{enumerate}

\subsection{From-Scratch Training}

Unlike the original DeepEF which used a pretrained checkpoint (43\_final\_model.pt) from native/decoy discrimination on CASP12, this work trains \textbf{entirely from scratch} on MegaScale mutation data. This means:
\begin{itemize}
    \item No transfer learning from structure discrimination
    \item Model must learn both structural representations AND stability prediction simultaneously
    \item Expected performance ceiling lower than with pretraining ($\sim$0.08--0.12 PCC gap)
\end{itemize}

\section{Data Processing}

\subsection{MegaScale/PNAS Dataset}

The dataset is derived from Tsuboyama et al.\ (2023) with the following filtering:
\begin{itemize}
    \item Remove insertions and deletions (point mutations only)
    \item Remove multi-site mutations (single-point only)
    \item Remove test set homologs from training
    \item Clamp $\Delta\Delta G$ to $[-1.0, 5.0]$ kcal/mol
    \item PNAS mutation list filtering for standardized comparison
\end{itemize}

Final dataset statistics:
\begin{itemize}
    \item \textbf{Training:} 340 proteins, $\sim$50,000--100,000 mutations
    \item \textbf{Test:} 28 proteins (from ThermoMPNN benchmark set)
    \item \textbf{Protein sizes:} 50--200 residues (small domains)
    \item \textbf{Structures:} AlphaFold2-predicted (available on Zenodo)
\end{itemize}

\subsection{Embedding Generation}

Currently: ProtT5 per-residue embeddings (1024-dim) for each mutant sequence, stored as:
\begin{verbatim}
training_data/PROTEIN_NAME/prott5_embeddings/
    prott5_embedding_0.pt   # Chunk 0: [128, seq_len, 1024]
    prott5_embedding_1.pt   # Chunk 1: [128, seq_len, 1024]
    ...
\end{verbatim}

Planned upgrade: SaProt structure-aware embeddings (1280-dim), requiring:
\begin{enumerate}
    \item Download AlphaFold PDB files from Zenodo (14 MB)
    \item Extract 3Di structural tokens via Foldseek
    \item Generate SaProt embeddings with combined AA+3Di input
\end{enumerate}
